%-------------------------
% Rover Resume 
% Link: https://github.com/subidit/rover-resume
%-------------------------

\usepackage[semibold]{raleway}
%\usepackage[default,semibold]{raleway}
\usepackage{fontspec}
\usepackage{fontawesome5}
\IfFileExists{../font/LinotypeSyntaxCom-Regular.ttf}{
  \setmainfont{LinotypeSyntaxCom}[
    Path=../font/,
    Extension = .ttf,
    UprightFont=*-Regular,
    BoldFont=*-Bold,
    ItalicFont=*-Italic,
    BoldItalicFont=*-BoldIt
  ]
}{
  \setmainfont{TeX Gyre Termes}
}

\usepackage[margin=1in,a4paper]{geometry}
\usepackage{enumitem}
\setlist[description]{itemsep=0pt,labelindent=1em}

\usepackage{xcolor}
\usepackage{titlesec}

\setcounter{secnumdepth}{0}

\titlespacing{\section}{0pt}{*4}{*0}
\titlespacing{\subsection}{0pt}{*3}{-\parskip}
\titlespacing{\subsubsection}{0pt}{\parskip}{-\parskip}
\titleformat*{\section}{\color{gray}\huge\fontseries{l}\scshape}
\titleformat*{\subsection}{\large\bfseries}
\titleformat*{\subsubsection}{\bfseries\scshape}

\newcommand{\aux}[1]{%
  $|$ {\normalfont \textsc{#1}}
}
\newcommand{\rside}[1]{
  \hfill {\normalfont \color{gray} \scshape\oldstylenums{\lowercase{#1}}}
}

\usepackage[hidelinks,
	pdfdisplaydoctitle = true,
	bookmarks=false,
	pdftitle={\title},
    pdfauthor={\Vorname~\Nachname}
	]{hyperref}

\usepackage{csquotes}
\usepackage{orcidlink}
\usepackage{polyglossia} % Language handling

\usepackage[
	style = ext-numeric-comp,
	backend=biber,
	sorting=ynt,
	date = year,
	articlein = false, %Suppress "In:" biblatex
	giveninits = true, % abbrevation for first name
	maxbibnames=99,
	minbibnames=99,
	maxcitenames=99,
	mincitenames=99,
	defernumbers=true]{biblatex}
%\AtBeginBibliography{\small}

\DeclareFieldFormat{title}{\mkbibemph{#1}}
\DeclareFieldFormat[article,incollection,inproceedings]{title}{\mkbibemph{#1}}
\DeclareFieldFormat{booktitle}{\mkbibemph{#1}}
\DeclareFieldFormat{journaltitle}{\mkbibemph{#1}}

% Allow name-based highlighting (robust substring matching)
\usepackage{etoolbox}
\usepackage{xstring}
\usepackage{hyphenat}
% replace - with hyphen in journal title to avoid line break issues
\DeclareSourcemap{
  \maps[datatype = bibtex]{
    \map{
     \step[fieldsource = journaltitle, match = \regexp{-}, replace = \regexp{\{\\hyp\{\}\}}]
    }
  }
}
% Ensure 'Sichtbeton' can be hyphenated in dynamically generated .bib entries
\DeclareSourcemap{
	\maps[datatype = bibtex]{
		\map{
			\step[fieldsource = title, match = \regexp{Sichtbeton}, replace = \regexp{Sicht\\-beton}]
		}
	}
}
\DeclareNameFormat{author}{%
	\nameparts{#1}%
	\edef\givenname{\namepartgiven}%
	\edef\familyname{\namepartfamily}%
	% If this name matches the full user name, print full given name and family name and underline
	\IfSubStr{\givenname}{\Vorname}{%
		\IfSubStr{\familyname}{\Nachname}{\underline{\StrLeft{\namepartgiven}{1}[\givenninit]\givenninit.\space\namepartfamily}}{%
			% Given name contains user's first name but family doesn't match: abbreviate given name
			\StrLeft{\namepartgiven}{1}[\givenninit]\givenninit.\space\namepartfamily}{}
	}{%
		% Otherwise abbreviate given name to initial + period
		\StrLeft{\namepartgiven}{1}[\givenninit]\givenninit.\space\namepartfamily
	}%
	\ifthenelse{\value{listcount}<\value{liststop}}{\addcomma\space}{}%
}

% Use simple title emphasis (sanitization disabled - fix characters in .bib if needed)
\DeclareFieldFormat{title}{\mkbibemph{#1}}
\DeclareFieldFormat[article,incollection,inproceedings]{title}{\mkbibemph{#1}}
\DeclareFieldFormat{booktitle}{\mkbibemph{#1}}
\DeclareFieldFormat{journaltitle}{\mkbibemph{#1}}

% remove ISBN from printed bibliography
\AtEveryBibitem{\clearfield{isbn}}

% https://tex.stackexchange.com/questions/134191/line-breaks-of-long-urls-in-biblatex-bibliography
\setcounter{biburllcpenalty}{7000}
\setcounter{biburlucpenalty}{8000}


% Footer: Lebenslauf info + page x/y in 80% gray
\usepackage{fancyhdr}
\usepackage{lastpage}
\pagestyle{fancy}
\fancyhf{}%
\renewcommand{\headrulewidth}{0pt}
\renewcommand{\footrulewidth}{0pt}
